\chapter{Méthodologie de l'épreuve}

\noindent\textit{Bien connaître son cours ne suffit pas : réussir l'épreuve de
Sciences de la Vie et de la Terre, de l'Éducation à l'Environnement et à la Santé du
BEPC demande aussi de comprendre comment elle est construite, ce que chaque mot de
consigne exige, et comment gérer son temps. C'est l'objet de ce chapitre.}
\vspace{8pt}

\section{Structure de l'épreuve}

\begin{propriete}[Ce qu'il faut savoir]
L'épreuve de SVTEEHB du BEPC est notée sur \textbf{20 points}. Elle comporte deux
parties \textbf{indépendantes et obligatoires} :
\begin{itemize}
\item \textbf{Partie A — Évaluation des ressources} (10 points), elle-même divisée en
deux volets : \textbf{I. Connaissances} (4 points, questions de cours : définitions,
schémas, propriétés) et \textbf{II. Compétences méthodologiques} (6 points, exploitation
de documents : tableaux, graphiques, schémas à analyser).
\item \textbf{Partie B — Évaluation des compétences} (10 points) : une
\textbf{situation-problème} ancrée dans un contexte concret (santé, environnement),
avec plusieurs \textbf{tâches} à accomplir, qui demande de \textbf{mobiliser plusieurs
notions ensemble} pour répondre à un problème de la vie courante.
\end{itemize}
\end{propriete}

\begin{remarque}
Les deux parties sont \textbf{indépendantes} : un échec sur une question de la Partie
A n'empêche pas de réussir la Partie B, et inversement. Ne jamais abandonner une
partie à cause d'une difficulté rencontrée dans l'autre.
\end{remarque}

\section{Lexique des consignes}

\begin{propriete}[Ce que chaque mot exige vraiment]
Le barème sanctionne autant la \textbf{qualité de la réponse au mot employé} que le
contenu scientifique. Voici ce qu'attend chaque consigne :
\end{propriete}

\begin{center}
\small
\begin{tabular}{@{}p{3.1cm}p{9.2cm}@{}}
\toprule
\textbf{Consigne} & \textbf{Ce qui est attendu} \\
\midrule
Nommer / Citer & Un mot ou une expression précise, sans phrase développée. \\
Définir & La définition exacte du terme, avec le vocabulaire scientifique correct. \\
Décrire & Une description ordonnée et complète d'une structure ou d'un
phénomène. \\
Expliquer / Justifier & Une phrase précise expliquant \textbf{pourquoi}, en
s'appuyant sur un mécanisme du cours. \\
Comparer & Mettre en relation deux éléments (avec ressemblances et différences),
et conclure. \\
Schématiser / Légender & Un schéma clair, proportionné et correctement légendé. \\
Interpréter un document & Relier les données du document (tableau, graphique) à une
connaissance du cours pour en tirer une conclusion. \\
Calculer & Un résultat numérique, avec la formule utilisée et le détail du
calcul. \\
Proposer & Une solution argumentée et réaliste, adaptée au contexte donné. \\
Conclure & Une phrase de synthèse répondant explicitement à la question posée. \\
\bottomrule
\end{tabular}
\end{center}

\begin{remarque}
« Interpréter un document » n'attend jamais une simple description de ce qui est vu :
il faut toujours \textbf{relier l'observation à une explication scientifique} tirée du
cours.
\end{remarque}

\section{Gérer son temps}

\begin{propriete}[Répartition indicative]
\begin{itemize}
\item \textbf{5 min} : lecture complète du sujet, repérage des parties les plus
faciles.
\item \textbf{20 min} : Partie A, I. Connaissances (questions de cours).
\item \textbf{35 min} : Partie A, II. Compétences méthodologiques (exploitation de
documents).
\item \textbf{50 min} : Partie B (situation-problème), en lisant \textbf{deux fois}
l'énoncé avant de commencer (les données utiles sont souvent réparties sur plusieurs
paragraphes ou documents).
\item \textbf{10 min} : relecture, vérification du vocabulaire scientifique et de la
présentation.
\end{itemize}
\end{propriete}

\begin{remarque}
Sur une question bloquante, ne pas s'acharner plus de $3$ à $4$ minutes : passer à la
suite, et revenir dessus s'il reste du temps. Une question non traitée fait perdre ses
points ; s'obstiner dessus fait perdre aussi ceux des questions suivantes, non lues.
\end{remarque}

\section{Rédiger et présenter sa copie}

\begin{propriete}[Bonnes pratiques]
\begin{itemize}
\item Utiliser le \textbf{vocabulaire scientifique précis} du cours (« phagocytose »,
« hématose », « diagenèse »\dots) plutôt que des formulations approximatives.
\item Pour un schéma, toujours \textbf{légender avec des traits de rappel} propres et
un titre.
\item Ne jamais laisser une question \textbf{blanche} : même une réponse partielle
peut rapporter un point.
\item Dans la situation-problème (Partie B), \textbf{revenir au contexte} : conclure
par une phrase reliée à la question concrète posée, pas seulement une liste de faits.
\item Structurer les réponses longues en phrases complètes, avec une
\textbf{conclusion explicite}.
\end{itemize}
\end{propriete}

\section{Dix pièges fréquents à éviter}

\begin{synthese}[Les erreurs qui coûtent le plus de points]
\textbf{1.} Confondre caractère héréditaire (transmis par les gènes) et caractère
acquis (dû à l'environnement) (chapitre 1).\\[3pt]
\textbf{2.} Croire qu'un caryotype à 47 chromosomes signifie qu'il manque un
chromosome, alors que la trisomie 21 correspond à un chromosome \textbf{en trop}
(chapitre 2).\\[3pt]
\textbf{3.} Oublier que deux parents porteurs sains (AS) peuvent avoir des enfants
sains (AA) et pas seulement des enfants malades (SS) (chapitre 3).\\[3pt]
\textbf{4.} Confondre la mitose (cellules filles identiques) et la formation des
gamètes avec brassage chromosomique (cellules différentes) (chapitre 4).\\[3pt]
\textbf{5.} Confondre contamination (entrée du microbe) et infection (développement
avec symptômes) (chapitre 5).\\[3pt]
\textbf{6.} Confondre asepsie (préventive, avant contact) et antisepsie (curative,
après contact) (chapitre 6).\\[3pt]
\textbf{7.} Confondre réponse non spécifique (immédiate, identique pour tout microbe)
et réponse spécifique (ciblée, avec mémoire immunitaire) (chapitre 7).\\[3pt]
\textbf{8.} Confondre sérothérapie (anticorps injectés, immédiat, courte durée) et
vaccination (anticorps fabriqués par l'organisme, différé, longue durée)
(chapitre 10).\\[3pt]
\textbf{9.} Confondre infarctus du myocarde (touche le cœur) et AVC (touche le
cerveau) (chapitre 11).\\[3pt]
\textbf{10.} Croire que le paludisme et la fièvre Ebola se transmettent de la même
façon : le paludisme par un vecteur (moustique), Ebola par contact direct
(chapitre 14).
\end{synthese}

\begin{remarque}
Chacun de ces pièges est repris, avec un exemple corrigé, dans l'encadré « Piège
fréquent » de la fiche de révision du chapitre concerné.
\end{remarque}
